\section{Proof of the Main Theorem}

\begin{proof}[Proof of Theorem 2.1]
Theorem \ref{thm:classical-closure} is the route-closing implication attached to the four source-grounded bridge propositions. Proposition \ref{prop:scale-barrier} suppresses the dangerous high-frequency cascade on the classical equation itself. Proposition \ref{prop:monotone-functional} keeps the global energy/enstrophy layer explicit on that same surface through the bounded control quantity \(Q(t)\) and hands the route forward to the post-\(Q(t)\) cascade carrier. Proposition \ref{prop:compactness} closes the quadratic nonlinear term on the same classical approximation surface, so no defect survives in the convective term. Proposition \ref{prop:gradient-control} records the Euclidean gradient-control fourth bridge and its role inside the original loop.

Accordingly, if those four bridge propositions are all available on one
admissibility surface with no hidden modified-equation work and no silent
branch substitution, then the only singularity route tracked in the present
framework is blocked on the declared theorem surface. That is exactly the
conditional implication asserted in Theorem \ref{thm:classical-closure}.
\end{proof}

\begin{remark}[Proof architecture]
The proof architecture is a four-way loop, not a linear pipeline. The scale barrier supplies the dangerous-scale suppression, the monotone functional keeps the global energy/enstrophy layer explicit, the compactness package closes the nonlinear term on the same equation, and the Euclidean gradient-control bridge controls gradient growth and feeds back into the scale barrier.
\end{remark}
